% ======================================================================
%  Protocolos de comunicación de bajo y alto nivel
%  EL5841 - Taller de Sistemas Embebidos, ITCR
%  Compilar con: pdflatex (dos pasadas) o latexmk -pdf protocolos.tex
% ======================================================================
\documentclass[aspectratio=169,11pt]{beamer}

\usetheme{metropolis}
\usepackage[spanish]{babel}
\usepackage[utf8]{inputenc}
\usepackage[T1]{fontenc}
\usepackage{booktabs}
\usepackage{multirow}
\usepackage{tikz}
\usepackage{listings}
\usepackage{xcolor}
\usetikzlibrary{arrows.meta,positioning,shapes.geometric,fit,backgrounds,decorations.pathreplacing}

% ---- Paleta ----------------------------------------------------------
\definecolor{fondo}{HTML}{FAFAF7}
\definecolor{tinta}{HTML}{1B3A2F}      % verde profundo
\definecolor{cobre}{HTML}{B0721A}      % acento cobre
\definecolor{acero}{HTML}{40607A}      % azul acero (alto nivel)
\definecolor{linea}{HTML}{9DB3A5}
\definecolor{suave}{HTML}{EBF0EA}

\setbeamercolor{normal text}{fg=tinta,bg=fondo}
\setbeamercolor{alerted text}{fg=cobre}
\setbeamercolor{frametitle}{bg=tinta,fg=fondo}
\setbeamercolor{progress bar}{fg=cobre}
\setbeamercolor{title separator}{fg=cobre}
\metroset{block=fill,progressbar=frametitle}

% ---- Código ----------------------------------------------------------
\lstset{
  basicstyle=\ttfamily\small,
  keywordstyle=\color{cobre}\bfseries,
  commentstyle=\color{acero}\itshape,
  stringstyle=\color{acero},
  showstringspaces=false,
  breaklines=true,
  frame=none,
  backgroundcolor=\color{suave},
  xleftmargin=0.6em,xrightmargin=0.6em,
  aboveskip=0.6em,belowskip=0.4em,
}

% ---- Ayudas visuales -------------------------------------------------
\newcommand{\capa}[2]{\textbf{\color{cobre}#1}\, #2}
\newcommand{\complejo}[1]{{\color{cobre}\textbf{Punto complejo:}} #1}
\newcommand{\curso}[1]{\begin{block}{En el curso}#1\end{block}}

\title{Protocolos de comunicación}
\subtitle{De los cables a la red: de bajo a alto nivel}
\author{Luis Diego García Rojas}
\institute{EL5841 — Taller de Sistemas Embebidos \\ Instituto Tecnológico de Costa Rica}
\date{II Semestre 2026}

\begin{document}

\maketitle

% ======================================================================
\begin{frame}{El sistema que nos va a acompañar}
Un \alert{sistema de control de acceso} sobre Raspberry Pi 4 con Linux embebido.

\vspace{0.6em}
\begin{center}
\begin{tikzpicture}[font=\small,node distance=6mm,
  chip/.style={draw=tinta,fill=suave,rounded corners,minimum height=8mm,inner sep=4pt},
  pi/.style={draw=tinta,fill=tinta,text=fondo,rounded corners,minimum height=14mm,minimum width=26mm,font=\bfseries}]
  \node[pi] (pi) {Raspberry Pi};
  \node[chip,above left=5mm and 14mm of pi] (sen) {Sensor};
  \node[chip,left=14mm of pi] (mem) {Flash};
  \node[chip,below left=5mm and 14mm of pi] (con) {Consola};
  \node[chip,above right=6mm and 16mm of pi] (nube) {Broker / Nube};
  \node[chip,below right=6mm and 16mm of pi] (tec) {Técnico};

  \draw[-{Latex}] (sen) -- (pi) node[midway,above,sloped,font=\tiny\bfseries,text=cobre]{I2C};
  \draw[-{Latex}] (mem) -- (pi) node[midway,above,font=\tiny\bfseries,text=cobre]{SPI};
  \draw[-{Latex}] (con) -- (pi) node[midway,below,sloped,font=\tiny\bfseries,text=cobre]{UART};
  \draw[{Latex}-{Latex}] (pi) -- (nube) node[midway,above,sloped,font=\tiny\bfseries,text=acero]{MQTT};
  \draw[{Latex}-{Latex}] (pi) -- (tec) node[midway,below,sloped,font=\tiny\bfseries,text=acero]{SSH};
\end{tikzpicture}
\end{center}

\vspace{0.3em}
{\footnotesize Cada protocolo de hoy resuelve una de estas conexiones. Al final volvemos a este diagrama.}
\end{frame}

% ======================================================================
\begin{frame}{Ruta de la charla}
\begin{columns}[T]
\column{0.5\textwidth}
\textbf{\color{cobre}Bajo nivel} \\ \emph{entre chips, capas 1--2}
\begin{itemize}
  \item UART / USART
  \item SPI
  \item I2C
  \item JTAG (interfaz de prueba y depuración)
\end{itemize}
\column{0.5\textwidth}
\textbf{\color{acero}Alto nivel} \\ \emph{entre sistemas, capa 7}
\begin{itemize}
  \item SSH
  \item MQTT
\end{itemize}
\end{columns}

\vspace{1em}
\begin{alertblock}{Hilo conductor}
Subimos por las capas del modelo OSI \emph{y} seguimos la vida del dispositivo: desde el primer arranque hasta la operación en campo.
\end{alertblock}
\end{frame}

% ======================================================================
\section{El modelo OSI}

\begin{frame}{¿Qué es un protocolo? ¿Por qué en capas?}
\begin{columns}[T]
\column{0.54\textwidth}
Un \alert{protocolo} es un conjunto de reglas acordadas para comunicarse.

\vspace{0.5em}
Se organizan en \alert{capas}: cada una resuelve un problema y confía en la de abajo, sin necesitar saber cómo funciona.

\vspace{0.5em}
{\footnotesize Cambiar el WiFi por un cable no obliga a reescribir la aplicación: solo cambia la capa de abajo.}

\column{0.46\textwidth}
\begin{block}{Idea clave}
Los protocolos de \textbf{bajo nivel} mueven \emph{bits entre chips}.

Los de \textbf{alto nivel} mueven \emph{mensajes entre sistemas}, apoyados en TCP/IP.
\end{block}
\end{columns}
\end{frame}

\begin{frame}{Dónde vive cada protocolo}
\begin{center}
\begin{tabular}{@{}ll@{}}
\toprule
\textbf{Capa OSI} & \textbf{Protocolos de hoy} \\
\midrule
7 — Aplicación & {\color{acero}\textbf{SSH, MQTT}} \\
4 — Transporte & TCP \emph{(base de ambos)} \\
3 — Red & IP \\
2 — Enlace & \multirow{2}{*}{{\color{cobre}\textbf{UART, SPI, I2C}}\; · \; JTAG$^{*}$} \\
1 — Física & \\
\bottomrule
\end{tabular}
\end{center}

\vspace{0.6em}
\begin{alertblock}{$^{*}$ Aclaración honesta}
JTAG \textbf{no} es un protocolo de datos como los demás: es una interfaz de \emph{prueba y depuración}. Lo incluimos porque es esencial en el ciclo de vida del sistema.
\end{alertblock}

\vspace{0.3em}
{\footnotesize Este mapa es el índice: vamos a volver a él al iniciar cada sección.}
\end{frame}

\begin{frame}{OSI y el stack de Linux embebido}
\begin{center}
\begin{tikzpicture}[font=\footnotesize,every node/.style={minimum width=6.4cm,minimum height=8mm,draw=tinta,anchor=west}]
  \node[fill=acero!18] (app) at (0,3.0) {\textbf{Aplicaciones} (espacio de usuario): SSH, MQTT};
  \node[fill=suave]     (ker) at (0,2.0) {\textbf{Kernel}: drivers, pila TCP/IP, \texttt{/dev/i2c-1}};
  \node[fill=cobre!18]  (dt)  at (0,1.0) {\textbf{Device tree}: declara los buses de la placa};
  \node[fill=linea!35]  (hw)  at (0,0.0) {\textbf{Hardware}: pines, UART, SPI, I2C, JTAG};
  \draw[{Latex}-{Latex},cobre,thick] (7.0,0.3) -- (7.0,2.7) node[midway,right,text=tinta]{sube};
\end{tikzpicture}
\end{center}
\vspace{0.4em}
{\footnotesize Los protocolos de bajo nivel se tocan desde \emph{drivers} o nodos en \texttt{/dev}; SSH y MQTT, desde \emph{aplicaciones}.}
\end{frame}

% ======================================================================
\section{Bajo nivel: comunicación entre chips}

\begin{frame}{Dos ejes que ordenan todo el bloque}
\begin{columns}[T]
\column{0.5\textwidth}
\begin{block}{¿Hay reloj compartido?}
\textbf{Asíncrono} (UART): no. \\
\textbf{Síncrono} (SPI, I2C): sí.
\end{block}
\column{0.5\textwidth}
\begin{block}{¿Punto a punto o bus?}
\textbf{Punto a punto} (UART, SPI). \\
\textbf{Bus} (I2C): muchos en dos líneas.
\end{block}
\end{columns}
\vspace{0.8em}
\begin{center}
{\color{cobre}\textbf{UART}} $\;\longrightarrow\;$ {\color{cobre}\textbf{SPI}} $\;\longrightarrow\;$ {\color{cobre}\textbf{I2C}} \\[0.3em]
{\footnotesize de lo más simple a lo más complejo}
\end{center}
\end{frame}

% ---- UART ------------------------------------------------------------
\begin{frame}{UART / USART}
\begin{columns}[T]
\column{0.55\textwidth}
\textbf{Fundamentos}
\begin{itemize}\small
  \item Asíncrono, sin reloj. Dos líneas: TX y RX.
  \item Trama: bit de inicio, datos, paridad opcional, bit de parada.
  \item Ambos lados acuerdan el \alert{baud rate}.
  \item USART agrega un modo síncrono.
\end{itemize}
\column{0.45\textwidth}
\textbf{Ejemplos y casos}
\begin{itemize}\small
  \item Consola de depuración, GPS, módulos Bluetooth.
  \item \textbf{Industrial:} UART sobre RS-485 es la base de \emph{Modbus RTU} en PLCs y variadores.
\end{itemize}
\end{columns}
\vspace{0.4em}
\complejo{sin reloj, el receptor detecta el flanco del bit de inicio y muestrea a mitad de cada bit (sobremuestreo $\times$16). De ahí la tolerancia de $\pm$2--3\,\% en el baud rate.}
\end{frame}

\begin{frame}{UART: la trama en el tiempo}
\begin{center}
\begin{tikzpicture}[font=\footnotesize]
  \def\u{0.85}
  \draw[->] (0,0) -- (12.6,0) node[right]{$t$};
  \draw[thick,cobre]
    (0,1) -- (1,1)                       % reposo (alto)
    -- (1,0) -- (2,0)                    % START (bajo)
    -- (2,1) -- (3,1)                    % b0=1
    -- (3,0) -- (5,0)                    % b1,b2=0
    -- (5,1) -- (6,1)                    % b3=1
    -- (6,0) -- (9,0)                    % b4..b6=0
    -- (9,1) -- (10,1)                   % b7=1
    -- (10,1) -- (11,1)                  % STOP (alto)
    -- (11,1) -- (12.3,1);               % reposo
  \node[below,font=\tiny] at (1.5,0)  {START};
  \foreach \x/\b in {2.5/1,3.5/0,4.5/0,5.5/1,6.5/0,7.5/0,8.5/0,9.5/1}
     \node[font=\tiny,text=tinta] at (\x,1.25) {\b};
  \node[above,font=\tiny] at (10.5,1) {STOP};
  \draw[decorate,decoration={brace,amplitude=4pt},acero] (2,-0.55) -- (10,-0.55)
     node[midway,below,font=\tiny]{8 bits de datos (LSB primero)};
\end{tikzpicture}
\end{center}
\vspace{0.3em}
{\footnotesize Se transmite primero el bit menos significativo. La línea en reposo está en alto.}
\end{frame}

% ---- SPI -------------------------------------------------------------
\begin{frame}{SPI}
\begin{columns}[T]
\column{0.55\textwidth}
\textbf{Fundamentos}
\begin{itemize}\small
  \item Síncrono y full-duplex.
  \item Líneas: SCLK, MOSI, MISO y un CS por periférico.
  \item Un controlador, varios periféricos. Muy rápido (decenas de MHz).
\end{itemize}
\column{0.45\textwidth}
\textbf{Ejemplos y casos}
\begin{itemize}\small
  \item Memorias flash, tarjetas SD, pantallas, ADC de alta velocidad.
\end{itemize}
\end{columns}
\vspace{0.4em}
\complejo{los \alert{4 modos (CPOL/CPHA)} definen en qué nivel reposa el reloj y en qué flanco se captura el dato. Es la causa \#1 de ``conecté todo y no funciona''.}
\end{frame}

\begin{frame}{SPI: el reloj sincroniza cada bit}
\begin{center}
\begin{tikzpicture}[font=\footnotesize]
  % reloj
  \node[left] at (0,2) {SCLK};
  \draw[thick,acero] (0,1.7)
    \foreach \x in {0,1,2,3,4,5,6,7} { -- (\x+0.5,1.7) -- (\x+0.5,2.3) -- (\x+1,2.3) -- (\x+1,1.7) };
  % datos
  \node[left] at (0,0.7) {MOSI};
  \draw[thick,cobre] (0,0.4)
    -- (1,0.4) -- (1,1.0) -- (3,1.0) -- (3,0.4) -- (5,0.4)
    -- (5,1.0) -- (6,1.0) -- (6,0.4) -- (8,0.4);
  \foreach \x/\b in {1.5/1,2.5/1,3.5/0,4.5/0,5.5/1,6.5/0,7.5/0}
     \node[font=\tiny] at (\x,0.7) {\b};
  \draw[dashed,linea] (0.5,0.2) -- (0.5,2.4);
  \node[font=\tiny,text=tinta] at (0.9,2.55) {captura en el flanco};
\end{tikzpicture}
\end{center}
\vspace{0.3em}
{\footnotesize El receptor no adivina el ritmo: el reloj le dice exactamente cuándo leer.}
\end{frame}

% ---- I2C -------------------------------------------------------------
\begin{frame}{I2C}
\begin{columns}[T]
\column{0.55\textwidth}
\textbf{Fundamentos}
\begin{itemize}\small
  \item Dos líneas (SDA, SCL) compartidas por muchos dispositivos.
  \item Direccionamiento de 7 bits, ACK/NACK.
  \item 100\,kHz, 400\,kHz, 1\,MHz.
\end{itemize}
\column{0.45\textwidth}
\textbf{Ejemplos y casos}
\begin{itemize}\small
  \item Sensores de temperatura y aceleración, EEPROM, RTC, gestión de energía.
\end{itemize}
\end{columns}
\vspace{0.4em}
\complejo{lógica \alert{open-drain con resistencias pull-up}: nadie ``empuja'' un 1, solo suelta la línea. Habilita el arbitraje multi-master y el \emph{clock stretching}, pero la capacitancia del bus limita la velocidad.}
\end{frame}

\begin{frame}{Comparativa del bloque de bajo nivel}
\begin{center}
\begin{tabular}{@{}lccc@{}}
\toprule
 & \textbf{UART} & \textbf{SPI} & \textbf{I2C} \\
\midrule
Reloj        & Asíncrono & Síncrono & Síncrono \\
Líneas       & 2 & 3 + 1/disp. & 2 \\
Dispositivos & 2 & Varios (CS) & Muchos (dirección) \\
Velocidad    & Baja--media & Muy alta & Media \\
Dúplex       & Full & Full & Half \\
Uso típico   & Consola, GPS & Flash, pantalla & Sensores \\
\bottomrule
\end{tabular}
\end{center}
\vspace{0.5em}
{\footnotesize Esta es, casi siempre, la diapositiva que más gente fotografía.}
\end{frame}

% ---- JTAG ------------------------------------------------------------
\begin{frame}{JTAG — ¿y cómo programamos y probamos el chip?}
\begin{columns}[T]
\column{0.55\textwidth}
\textbf{Fundamentos}
\begin{itemize}\small
  \item Estándar IEEE 1149.1. Líneas TCK, TMS, TDI, TDO.
  \item Nació para probar placas (\emph{boundary scan}); hoy también programa y depura.
  \item SWD: alternativa de 2 pines en ARM.
\end{itemize}
\column{0.45\textwidth}
\textbf{Ejemplos y casos}
\begin{itemize}\small
  \item Prueba de soldaduras en producción sin sondas físicas.
  \item Programación de FPGA, depuración paso a paso.
\end{itemize}
\end{columns}
\vspace{0.4em}
\complejo{la \alert{máquina de estados del TAP} (16 estados manejados solo con TMS) y el encadenamiento de varios chips (\emph{daisy chain}).}
\vspace{0.3em}
\begin{alertblock}{Puente a seguridad}
Un JTAG habilitado en un producto comercial permite \emph{extraer el firmware}: es un vector de ataque real. Nos lleva de la mano a SSH.
\end{alertblock}
\end{frame}

% ======================================================================
\section{Transición: de los cables a la red}

\begin{frame}{¿Por qué no alcanza con lo anterior?}
\begin{columns}[c]
\column{0.5\textwidth}
UART, SPI e I2C funcionan en \alert{centímetros}, entre chips de una placa.
\vspace{0.6em}

Para llegar a un servidor al otro lado del mundo hace falta:
\begin{itemize}\small
  \item Direccionamiento global — \textbf{IP}
  \item Entrega confiable — \textbf{TCP}
  \item Seguridad — \textbf{cifrado y autenticación}
\end{itemize}
\column{0.5\textwidth}
\begin{center}
\begin{tikzpicture}[font=\footnotesize,node distance=4mm]
  \node[draw=tinta,fill=suave,rounded corners,minimum width=2.4cm,minimum height=8mm] (chip){cm: chip a chip};
  \node[draw=acero,fill=acero!15,rounded corners,minimum width=2.4cm,minimum height=8mm,below=8mm of chip] (red){km: red global};
  \draw[-{Latex},cobre,thick] (chip) -- (red) node[midway,right]{TCP/IP};
\end{tikzpicture}
\end{center}
\end{columns}
\vspace{0.6em}
{\footnotesize Por eso SSH y MQTT viven \emph{sobre} TCP/IP, en la capa de aplicación.}
\end{frame}

% ======================================================================
\section{Alto nivel: comunicación entre sistemas}

% ---- SSH -------------------------------------------------------------
\begin{frame}{SSH}
\begin{columns}[T]
\column{0.55\textwidth}
\textbf{Fundamentos}
\begin{itemize}\small
  \item Cliente-servidor, puerto 22.
  \item Canal cifrado y autenticado sobre una red insegura.
  \item Reemplazó a Telnet (contraseñas en texto plano).
\end{itemize}
\column{0.45\textwidth}
\textbf{Ejemplos y casos}
\begin{itemize}\small
  \item Administración remota de gateways y Raspberry Pi.
  \item Actualización con SCP/SFTP, Git, túneles.
\end{itemize}
\end{columns}
\vspace{0.4em}
\complejo{acordar una clave secreta sin haberse visto (\alert{Diffie-Hellman}); verificar la identidad del servidor con su \emph{host key} (la ``huella'' de la primera conexión); autenticar por llave pública en vez de contraseña.}
\end{frame}

% ---- MQTT ------------------------------------------------------------
\begin{frame}{MQTT}
\begin{columns}[T]
\column{0.55\textwidth}
\textbf{Fundamentos}
\begin{itemize}\small
  \item Modelo \alert{publicar/suscribir} con un \textbf{broker} en el centro.
  \item Los dispositivos hablan por \emph{tópicos}: \texttt{planta1/motor3/vibracion}.
  \item Liviano, para redes inestables. Puerto 1883, o 8883 con TLS.
\end{itemize}
\column{0.45\textwidth}
\textbf{Ejemplos y casos}
\begin{itemize}\small
  \item Monitoreo industrial (IIoT), domótica.
  \item AWS IoT, Azure IoT, Mosquitto. Sparkplug B para SCADA.
\end{itemize}
\end{columns}
\vspace{0.4em}
\complejo{niveles de \alert{QoS} (0: a lo sumo una vez; 1: al menos una vez; 2: exactamente una vez), \emph{retained messages} y \emph{Last Will} (aviso de desconexión inesperada).}
\end{frame}

\begin{frame}{SSH frente a MQTT}
\begin{center}
\begin{tabular}{@{}lcc@{}}
\toprule
 & \textbf{SSH} & \textbf{MQTT} \\
\midrule
Modelo    & Uno a uno & Muchos a muchos \\
Propósito & \emph{Controlar} un equipo & \emph{Mover} datos \\
Patrón    & Sesión interactiva & Publicar / suscribir \\
Intermediario & No & Broker \\
\bottomrule
\end{tabular}
\end{center}
\vspace{0.6em}
{\footnotesize No compiten: resuelven problemas distintos de la misma capa.}
\end{frame}

% ======================================================================
\section{Cierre y actividad}

\begin{frame}{El recorrido completo de un evento}
\begin{center}
\begin{tikzpicture}[font=\footnotesize,node distance=8mm,
  b/.style={draw=tinta,fill=suave,rounded corners,minimum height=8mm,inner sep=4pt}]
  \node[b] (s) {Sensor};
  \node[b,right=of s] (pi) {Raspberry Pi};
  \node[b,right=of pi] (br) {Broker};
  \node[b,right=of br] (dash) {Dashboard};
  \draw[-{Latex}] (s)--(pi) node[midway,above,font=\tiny,text=cobre]{I2C};
  \draw[-{Latex}] (pi)--(br) node[midway,above,font=\tiny,text=acero]{MQTT/TCP/IP};
  \draw[-{Latex}] (br)--(dash) node[midway,above,font=\tiny,text=acero]{MQTT};
  \node[below=10mm of pi,text=tinta] (tec) {Técnico};
  \draw[{Latex}-{Latex}] (tec)--(pi) node[midway,right,font=\tiny,text=acero]{SSH};
\end{tikzpicture}
\end{center}
\vspace{0.6em}
\begin{alertblock}{Mensaje central}
No hay un protocolo ``mejor''. Cada uno resuelve el problema de su capa, y un sistema real los necesita todos.
\end{alertblock}
\end{frame}

\begin{frame}{La misión: laboratorio cerrado}
Un corte de luz bloqueó el laboratorio. La contraseña del servidor quedó \alert{dividida en dos}.
\vspace{0.5em}
\begin{enumerate}
  \item \textbf{I2C (Tinkercad):} el sensor de temperatura guarda el \textbf{fragmento 1}. Se libera al llegar a 30\,°C.
  \item \textbf{UART (Tinkercad):} programan paso a paso el lector y la cerradura. Al abrir, revela el \textbf{fragmento 2}.
  \item \textbf{SSH (en casa, con Docker):} construyen el servidor archivo por archivo, entran con los dos fragmentos, abren la puerta y descargan el acta.
\end{enumerate}
\vspace{0.5em}
{\footnotesize Cada protocolo abre una parte del camino: si uno falla, la puerta sigue cerrada.}
\end{frame}

\begin{frame}{Referencias}
\footnotesize
\begin{itemize}
  \item H. Zimmermann, ``OSI Reference Model---The ISO Model of Architecture for Open Systems Interconnection'', \emph{IEEE Transactions on Communications}, vol. 28, no. 4, pp. 425--432, 1980.
  \item NXP Semiconductors, \emph{UM10204: I2C-bus specification and user manual}, Rev. 7.0, 2021.
  \item IEEE, \emph{IEEE Std 1149.1: Standard for Test Access Port and Boundary-Scan Architecture}.
  \item T. Ylonen y C. Lonvick, \emph{RFC 4251: The Secure Shell (SSH) Protocol Architecture}, IETF, 2006.
  \item T. Ylonen y C. Lonvick, \emph{RFC 4253: The Secure Shell (SSH) Transport Layer Protocol}, IETF, 2006.
  \item OASIS, \emph{MQTT Version 5.0}, OASIS Standard, 2019.
  \item Microchip Technology, \emph{ATmega328P Datasheet} (USART, SPI y TWI).
  \item Analog Devices, \emph{TMP35/TMP36/TMP37 Low Voltage Temperature Sensors}, hoja de datos.
  \item Arduino, documentación de referencia de las librerías \texttt{Wire} y \texttt{SoftwareSerial}, arduino.cc.
  \item OpenBSD Project, páginas de manual de OpenSSH: \texttt{ssh(1)}, \texttt{sshd\_config(5)}, \texttt{scp(1)}.
\end{itemize}
\end{frame}

\begin{frame}[standout]
  ¿Preguntas?
  \vspace{1em}

  {\normalsize Del cable al servidor, capa por capa.}
\end{frame}

\end{document}
